\oldpage[II]{217}
\phantomsection
\addcontentsline{toc}{section}{Erratas y adiciones (Lista 1)}
\section*{Erratas y adiciones}
\begin{center}
\large (Lista 1)
\end{center}

\subsection*{Capítulo 0}

\paragraph{(0, 2.1.3)}
En la línea 15 contando desde abajo de la p.~21, reemplazar
«inter-sección» por «intersección».

\paragraph{(0, 2.1.6)}
En la línea 13 de la p.~22, reemplazar
$Z_i\cap\complement(\bigcup\limits_{j\ne i}Z_j)$ par
$\bigcup\limits_{j\ne i}Z_j$.

\paragraph{(0, 3.1.6)}
En la línea 17 contando desde abajo de la p.~25, reemplazar
$\Gamma(U,\sh{F})$ por $\Gamma(V,\sh{F})$.

\paragraph{(0, 4.1.1)}
En la línea 12 de la p.~36, reemplazar $\sh{G}\to\sh{F}$ por
$\sh{B}\to\sh{A}$.

\paragraph{(0, 4.2.1)}
En la línea 11 contando desde abajo de la p.~39, reemplazar «haces» por
«haz».

\paragraph{(0, 5.1.2)}
Después de la línea 14 de la p.~45, añadir:

Si $\sh{G}$ es un $\sh{O}_Y$-módulo generado por sus secciones sobre
$Y$, entonces $f^*(\sh{G})$ está generado por sus secciones sobre $X$,
pues $f^*$ es un funtor exacto por la derecha.

\paragraph{(0, 5.2.4)}
En la línea 18 contando desde abajo de la p.~46, reemplazar $f_U$ por
$f_U^*$.

\paragraph{(0, 5.3.4)}
Antes de la línea 16 contando desde abajo de la p.~47, añadir:

Si
$\sh{A}\to\sh{B}\to\sh{C}\to\sh{D}\to\sh{E}$
es una sucesión exacta de $\sh{O}_X$-módulos, y si
$\sh{A},\sh{B},\sh{D},\sh{E}$ son coherentes, entonces $\sh{C}$ es
coherente.

\paragraph{(0, 5.4.1)}
Después de la línea 2 de la p.~49, añadir:

Supongamos que $\sh{O}_X$ es \emph{coherente}, y sea $\sh{F}$ un
$\sh{O}_X$-módulo \emph{coherente}. Si, en un punto $x\in X$, $\sh{F}_x$ es
un $\sh{O}_x$-módulo \emph{libre de rango $n$}, existe una vecindad $U$ de $x$
tal que $\sh{F}|U$ es \emph{localmente libre de rango $n$}; en efecto,
$\sh{F}_x$ es entonces isomorfo a $\sh{O}_x^n$, y la afirmación se
deduce de (5.2.7).

\paragraph{(0, 6.6.2)}
Reemplazar las líneas 10 y 11 de la p.~59 por:

En efecto, esto se deduce de (6.4.1, \emph{d)}).

\paragraph{(0, 6.7.1)}
Antes de la línea 10 contando desde abajo de la p.~59, añadir:

Si $f$ es un morfismo plano, también se dice que $X$ es
\emph{plano sobre $Y$}, o \emph{$Y$-plano}.

\paragraph{(0, 7.2.9)}
En la línea 11 contando desde abajo de la p.~65, reemplazar $M_n$ por
$M_{n-1}$.

\subsection*{Capítulo I}

\paragraph{(I, 1.2.1)}
En la línea 18 de la p.~83, reemplazar $\boldsymbol{K}$ por $\boldsymbol{k}$ (dos veces).

\paragraph{(I, 1.2.7)}
En la línea 16 contando desde abajo de la p.~84, reemplazar $A$ por
$A'$ (dos veces) y $\mathfrak N$ por $\mathfrak N'$; en la línea 17 contando desde
abajo, reemplazar $X$ por $X'$.

